Dans les hôpitaux et cliniques en Tunisie, il n’existe actuellement \textbf{aucun système de collecte de données automatisé} en salle de réveil. Le fonctionnement repose sur :
\begin{itemize}
    \item \textbf{Un écran de surveillance centralisé} installé dans une salle administrative, permettant au personnel de suivre les signes vitaux de plusieurs patients à distance.
    \item \textbf{Une surveillance manuelle régulière} : Le personnel médical (infirmiers et médecins) doit se déplacer physiquement d’un patient à l’autre pour vérifier les constantes vitales, ce qui entraîne des \textbf{délais} dans la détection des situations critiques.
\end{itemize}